% Generated draft: requires mathematical and source review.
\documentclass[11pt]{article}
\usepackage[margin=25mm]{geometry}
\usepackage{fontspec}
\setmainfont{Latin Modern Roman}
\usepackage{amsmath,amssymb,mathtools,unicode-math}
\setmathfont{Latin Modern Math}
\usepackage{enumitem,needspace,xurl,hyperref}
\setlength{\parindent}{0pt}
\setlength{\parskip}{5pt}
\setlength{\emergencystretch}{3em}
\newcommand{\N}{\mathbb{N}}
\newcommand{\Z}{\mathbb{Z}}
\newcommand{\Q}{\mathbb{Q}}
\newcommand{\R}{\mathbb{R}}
\newcommand{\C}{\mathbb{C}}
\newcommand{\F}{\mathbb{F}}
\newcommand{\A}{\mathbb{A}}
\newcommand{\PP}{\mathbb{P}}
\newcommand{\E}{\mathbb{E}}
\newcommand{\eps}{\varepsilon}
\newcommand{\dd}{\,\mathrm{d}}
\begin{document}
\section*{UnsolvedMath ID: 10000028}
\textbf{Source problem code:} catalogue entry
\subsection{Definitions and mathematical statement}
Let \((M,\rho)\) be a uniform infinite planar triangulation (UIPT), conventionally a local weak limit of uniformly sampled finite rooted planar triangulations in a fixed triangulation class. Conditional on \(M\), let \((X_n)\) be simple random walk started at \(X_0=\rho\), and let \(d_{\mathrm{gr}}\) be graph distance. The source asks for displacement exponent \(1/4\).

One precise subpolynomial, distributional version is
\[
 \forall\varepsilon>0:\quad
 \mathbb P\left(n^{1/4-\varepsilon}\le
 d_{\mathrm{gr}}(\rho,X_n)\le n^{1/4+\varepsilon}\right)
 \longrightarrow1\qquad(n\to\infty).
\]
With \(\mathbb P\) interpreted as the joint map-and-walk law, this is an annealed assertion. A corresponding almost-every-map conditional-probability formulation is quenched and needs separate specification. The display is an explicit candidate interpretation, not a claim that the excerpt states this exact mode.

The supplied wording combines \(d_{\mathrm{gr}}(X_0,X_n)\asymp n^{1/4}\) with ``subpolynomial or constant factors''. These are different strengths: an exponent up to \(n^{o(1)}\), fixed multiplicative comparability, and a nondegenerate rescaled limit are not interchangeable. Pointwise displacement at time \(n\), maximal displacement up to \(n\), and an eventual almost-sure lower bound are also distinct statements. The source does not fix them, and zero displacement at return times cannot be ignored when choosing an almost-sure formulation.

The UIPT convention concerning loops, multiple edges and rooting is not given and should be matched to the exact theorem. The summary labels the exponent solved, attributes the sharp result to Gwynne--Miller and earlier polynomial bounds to Benjamini--Curien, and mentions UIPQ as related. No constant-factor or scaling-limit theorem is inferred from this exponent report. The historical conjecture and solved-status attribution are retained pending exact primary-theorem verification.
\subsection{Short English statement}
Show that random-walk graph displacement on the UIPT grows with exponent \(1/4\). The supplied summary reports this exponent proved, but the exact map model, probability mode and subpolynomial-versus-constant-factor statement need to be fixed.
\subsection{Sources}
\begin{itemize}
\item Local registry: \url{data/unsolvedmath/problems.json}, permanent UnsolvedMath ID \(10000028\), source code catalogue entry; supplied statement, background, and research summary, updated 2026-08-06.
\item Itai Benjamini, \emph{Coarse Geometry and Randomness}, Saint-Flour notes, Conjecture 6.4, PDF page 49: \url{https://www.wisdom.weizmann.ac.il/~itai/stflouraug24.pdf#page=49}. Preserved copy: \url{https://arquivo.pt/noFrame/replay/20201231041548id_/http://www.wisdom.weizmann.ac.il/~itai/stflouraug24.pdf}. Supplied extraction checked against rendered PDF, accessed 2026-07-29.
\item Gwynne--Miller, ``Random walk on random planar maps: spectral dimension, resistance and displacement'', arXiv:1711.00836, cited by the summary for the sharp exponent; Benjamini--Curien work is mentioned without a separate full citation.
\end{itemize}
\end{document}
